% ============================================================
%  Chapter 1 -- Finite-State Automata and Regular Languages
% ============================================================

\chapter{Finite-State Automata and Regular Languages}
\label{ch:fsa}

\aside{Automata are the first modelling
language of the course: states are the
possible configurations of a system, and
transitions are the steps it can take.}

Verification starts with a simple modelling decision.  We describe a
system by listing the situations it can be in, and the events that move
it from one situation to another.  This gives us a directed graph with a
finite set of nodes, but with enough structure to process input words.
That structure is a \textbf{finite-state automaton}.

This chapter builds the basic language.  We begin with alphabets and
words, define deterministic and non-deterministic finite automata, and
prove the central fact that non-determinism does not add expressive
power for finite words.

\constructionpipeline{1}

\aside{Why start with finite words? Real-world systems like operating systems or network protocols often run continuously, effectively generating \emph{infinite} words. We build the foundations on finite sequences first, but will later extend this theory to infinite words to model systems that do not terminate.}

% ============================================================
\section{Alphabets, Words, and Languages}
\label{sec:alphabets-words}
% ============================================================

Before an automaton can read anything, we need to say what its inputs
look like.  For a turnstile, the relevant observations might be
\texttt{coin} and \texttt{push}; for a tiny parity example, they might
be just $a$ and $b$.  In both cases we first fix the symbols that may
appear, and then we study finite sequences of those symbols.

\begin{definition}[Alphabet and word]
\label{def:alphabet-word}
An \textbf{alphabet} is a finite non-empty set of symbols, usually
denoted by $A$.

A \textbf{word} over $A$ is a finite sequence of symbols from
$A$.  The set of all finite words over $A$ is denoted by
$A^*$.
\end{definition}

Think of $A$ as the vocabulary available to the automaton.  If
$A=\{a,b\}$, then $aab$, $baba$, and $bb$ are words.  The special
word with no symbols is the \textbf{empty word}, written $\varepsilon$.
It has length $0$.

\begin{example}[A small alphabet]
\label{ex:small-alphabet}
Let $A=\{\texttt{coin},\texttt{push}\}$.  A turnstile can observe
words such as
\[
  \texttt{coin}\;\texttt{push}\;\texttt{push}\;\texttt{coin}.
\]
The symbols are events.  The word is the finite history of observed
events.
\end{example}

\begin{definition}[Language]
\label{def:finite-language}
A \textbf{language} over $A$ is any set of finite words:
\[
  L \subseteq A^*.
\]
\end{definition}

\aside{In practice, the languages we care about are usually infinite.  A
language is still just a set of finite words, but it is typically an
infinite subset of $A^*$.}

A language can be understood as a way of assigning a property to words:
some words satisfy the property, while others do not. For example, over
the alphabet \(A=\{a,b\}\), we can form words such as \(a\), \(b\),
\(ab\), \(ba\), and \(abba\). The property ``ends in \(a\)'' selects
exactly those words whose last symbol is \(a\). Thus, the corresponding
language is the set of all and only the words satisfying that property:
\[
L = \{ w \in A^* \mid w \text{ ends in } a \}.
\]

In verification, this viewpoint is useful because a word can represent a
behaviour, that is, a sequence of states, events, or actions. A language
can therefore describe the set of behaviours allowed by a model, or the
set of behaviours permitted by a specification. Saying that a behaviour
belongs to the language means that it satisfies the property described by
that language.

From the automata perspective, a language is something a machine can
\emph{accept}.  An automaton reads an input word and decides whether that
word belongs to the language.  The languages accepted by finite automata
are called \textbf{regular languages}.  Later in the chapter we will see
that deterministic and non-deterministic finite automata recognize the
same class of languages, so ``regular'' is the right name for the whole
finite-state setting.

\aside{Finite automata are only one presentation of regular behaviour.
Later we will see regular expressions as an equivalent syntax for
describing the same class of languages.}

\begin{definition}[Basic operations on words]
\label{def:word-ops}
If $u,v\in A^*$, their \textbf{concatenation} is the word $uv$
obtained by writing $v$ immediately after $u$.  The \textbf{length} of
$w$ is denoted $|w|$.
\end{definition}

\aside{Concatenation is not commutative in general: usually $uv \neq vu$.
It is also length-additive: $|uv|=|u|+|v|$.}

The empty word behaves like a neutral element:
\[
  \varepsilon w = w\varepsilon = w.
\]
This small fact is useful because many constructions need to handle the
case where no input has been consumed yet.

The same operation extends from words to languages.  If
$L_1,L_2\subseteq A^*$, then
\[
  L_1L_2=\{uv:u\in L_1,\ v\in L_2\}.
\]
There is a small but important distinction here:
\[
  L\{\varepsilon\}=L,
  \qquad
  L\emptyset=\emptyset.
\]
The singleton $\{\varepsilon\}$ contains one word, namely the empty
word.  The empty language $\emptyset$ contains no word at all.

\begin{definition}[Kleene star]
\label{def:kleene-star}
If $W \subseteq A^*$ is a language, then
\[
  W^* =
  \{w_1w_2\cdots w_k : k\geq 0,\; w_i\in W\text{ for all }i\}
\]
is the \textbf{Kleene closure} of $W$.
\end{definition}

In plain language, $W^*$ contains every word that can be built by
concatenating finitely many blocks from $W$.  The case $k=0$ gives
$\varepsilon$, so $\varepsilon\in W^*$ for every $W$.

It is often useful to name the finite powers explicitly:
\[
  W^0=\{\varepsilon\},
  \qquad
  W^{n+1}=WW^n,
  \qquad
  W^*=\bigcup_{n\geq 0}W^n,
  \qquad
  W^+=\bigcup_{n\geq 1}W^n.
\]
For example, if $W=\{a,b\}$ and we want to describe words containing
at least one symbol, then $W^+$ is the natural choice. Using $W^*$
would also include the empty word $\varepsilon$, which would have to be
treated as a special case.

\begin{example}[Words from the alphabet $\{a\}$]
\label{ex:unary-words}
Let $A=\{a\}$.  Then $A^*=\{\varepsilon,a,aa,aaa,\ldots\}$ 
is the set of all finite sequences of $a$'s.
$A^+$ on the other hand is $\{a,aa,aaa,\ldots\}$, which excludes the empty word.
\end{example}


% ============================================================
\section{Deterministic Finite Automata}
\label{sec:dfa}
% ============================================================

A deterministic automaton has no choices.  Once we know the current
state and the next input symbol, the next state is fixed.  Before
writing the general definition, let us build one concretely.

\begin{example}[Words with an even number of $a$'s]
\label{ex:even-a-dfa}
Over $A=\{a,b\}$, consider the language of words containing an
even number of occurrences of $a$.  To decide membership we only need
to track the parity of the number of $a$'s seen so far.  This suggests
two states,
\[
  Q=\{q_{\mathrm{even}},q_{\mathrm{odd}}\},
\]
with $a$ toggling between them and $b$ leaving the state unchanged.
The initial state is $q_{\mathrm{even}}$, and it is also the only
accepting state, because zero is even.
\end{example}

\begin{figure}[ht]
\centering
\begin{tikzpicture}[
  node distance=3.2cm,
  state/.style={circle, draw, thick, minimum size=11mm, fill=blue!10},
  accepting/.style={circle, draw, double, thick, minimum size=11mm, fill=teal!10},
  arr/.style={-Stealth, thick},
]
  \node[accepting] (e) {$q_e$};
  \node[state, right=of e] (o) {$q_o$};
  \draw[arr] ([xshift=-1.0cm]e.west) -- (e);
  \draw[arr, bend left=18] (e) to node[above] {$a$} (o);
  \draw[arr, bend left=18] (o) to node[below] {$a$} (e);
  \draw[arr] (e) edge[loop above] node {$b$} (e);
  \draw[arr] (o) edge[loop above] node {$b$} (o);
\end{tikzpicture}
\caption{A DFA for words with an even number of $a$'s.}
\label{fig:even-a-dfa}
\end{figure}

In \cref{fig:even-a-dfa}, the state records only the information that
matters for the property: parity.  It forgets the exact number of $a$'s,
because the language does not need that much detail.  The picture
already contains every ingredient of a deterministic automaton: a finite
set of states, an alphabet of inputs, a rule that gives the next state,
a designated starting state, and a set of accepting states.  We now
formalize them.

\begin{definition}[Deterministic finite automaton]
\label{def:dfa}
A \textbf{deterministic finite automaton} (DFA) is a tuple
\[
  \mathcal{A}=(Q,A,\delta,q_0,F)
\]
where:
\begin{itemize}[nosep]
  \item $Q$ is a finite set of states;
  \item $A$ is a finite alphabet;
  \item $\delta:Q\times A\to Q$ is the transition function;
  \item $q_0\in Q$ is the initial state;
  \item $F\subseteq Q$ is the set of accepting states.
\end{itemize}
\end{definition}

The tuple separates the static shape of the automaton from its
behaviour.  The set $Q$ tells us which configurations exist.  The
transition function $\delta$ tells us how the configuration evolves
while reading a word.  The accepting set $F$ tells us which final states
count as success.

\aside{Using a single initial state is mostly a convenience.  Variants
with several initial states can be converted into this form without
changing the recognized language.}

\begin{example}[Every $a$ is eventually followed by a $b$]
\label{ex:a-eventually-followed-by-b}
Over $A=\{a,b\}$, consider the language of all and only those words such
that each occurrence of $a$ is followed later by an occurrence of $b$:
\[
  L=\{w\in\{a,b\}^* :
      \text{for every position }i,\;
      w_i=a \Rightarrow \exists j>i \text{ with } w_j=b\}.
\]
The word $\varepsilon$ belongs to $L$, and so does every word made only
of $b$'s, because there is no occurrence of $a$ to check.  The word
$aaab$ also belongs to $L$: the final $b$ witnesses all three pending
$a$'s at once.  By contrast, $a$, $aba$, and $baa$ do not belong to
$L$, because their last $a$ has no later $b$.
\end{example}

\begin{figure}[ht]
\centering
\begin{tikzpicture}[
  node distance=3.2cm,
  state/.style={circle, draw, thick, minimum size=11mm, fill=blue!10},
  accepting/.style={circle, draw, double, thick, minimum size=11mm, fill=teal!10},
  arr/.style={-Stealth, thick},
]
  \node[accepting] (ok) {$q_{\mathrm{ok}}$};
  \node[state, right=of ok] (wait) {$q_{\mathrm{wait}}$};
  \draw[arr] ([xshift=-1.0cm]ok.west) -- (ok);
  \draw[arr] (ok) edge[loop above] node {$b$} (ok);
  \draw[arr, bend left=18] (ok) to node[above] {$a$} (wait);
  \draw[arr] (wait) edge[loop above] node {$a$} (wait);
  \draw[arr, bend left=18] (wait) to node[below] {$b$} (ok);
\end{tikzpicture}
\caption{A DFA that remembers only whether there is an unresolved $a$.}
\label{fig:a-eventually-followed-by-b}
\end{figure}

The state $q_{\mathrm{wait}}$ means ``some $a$ has been seen and no
later $b$ has appeared yet''.  The automaton does not count the $a$'s
and the $b$'s.  It remembers only one bit of information: whether there
is currently an unresolved $a$.  This is exactly the finite-memory point.

\aside{The contrast with balanced parentheses is important.  There,
each opening parenthesis must be matched by the corresponding later
closing parenthesis, possibly after nested openings.  That requires an
unbounded stack of pending openings, so the language is not regular; it
belongs to the world of pushdown automata.}

\begin{definition}[Extended transition function]
\label{def:extended-delta}
The transition function extends from symbols to words by defining
$\widehat{\delta}:Q\times A^*\to Q$ recursively:
\[
  \widehat{\delta}(q,\varepsilon)=q,
  \qquad
  \widehat{\delta}(q,wa)=\delta(\widehat{\delta}(q,w),a).
\]
\end{definition}

This simply says: to read the word $wa$, first read $w$, then process the
last symbol $a$ from the state reached after $w$.
The sequence of states visited while reading a word is called a
\textbf{computation} of the automaton.

\begin{definition}[Language of a DFA]
\label{def:dfa-language}
The language accepted by a DFA $\mathcal{A}=(Q,A,\delta,q_0,F)$ is
\[
  L(\mathcal{A})=\{w\in A^* : \widehat{\delta}(q_0,w)\in F\}.
\]
\end{definition}

So acceptance is a final-state test.  The automaton consumes the whole
word; if the unique resulting state is accepting, the word belongs to
the language.

\subsection{Completeness and Sink States}
\label{subsec:sink-state}

\aside{In a formal DFA, every pair $(q,a)$ has exactly one next state
because $\delta$ is a total function.  In drawings, however, authors
often omit some arrows.  The issue is then not determinism but how to
complete the picture.}

Sometimes automata are drawn with some transitions omitted. This is often
done to keep the picture simple, but formally a DFA must have a
\emph{total} transition function. This means that for every state \(q\)
and every input symbol \(a \in A\), the transition \(\delta(q,a)\)
must be defined.

An omitted transition should therefore be understood as a transition that
goes to a rejecting ``dead'' or ``sink'' state. Once the automaton enters
this sink state, it stays there forever, no matter which input symbol is
read. Thus, any partially drawn DFA can be made into a formally correct
DFA by adding a sink state and redirecting all missing transitions to it.

\begin{definition}[Sink completion]
\label{def:sink-completion}
Given a partial deterministic automaton, its \textbf{completion} adds a
fresh non-accepting sink state $\bot$.  Every missing transition goes to
$\bot$, and every symbol loops from $\bot$ back to $\bot$.
\end{definition}

A missing transition is therefore not a self-loop.  Replacing omitted
arrows by self-loops can change the language.  To read this definition,
ask what the automaton should do when the drawing gives no arrow for the
next input symbol.  The operational test is simple: every missing case is
sent to the rejecting state $\bot$, and once $\bot$ is reached, the word
cannot be accepted.

\begin{example}[Why omitted transitions are not self-loops]
\label{ex:missing-transition-not-self-loop}
Let $A=\{a,b\}$.  Consider the partially drawn automaton with two
states:
\[
  q_0 \xrightarrow{a} q_1,
  \qquad
  q_1\in F.
\]
No other transition is drawn.  With sink completion, this automaton
recognizes exactly the one-letter word $a$: any unexpected symbol, such
as the first $b$ in $ba$ or the final $b$ in $ab$, sends the computation
to $\bot$ and the word is rejected.

If instead we completed the missing transitions by self-loops, the
language would be different.  The word $ba$ would be accepted, because
the missing $b$-transition from $q_0$ would keep the automaton in $q_0$,
and the following $a$ would move it to $q_1$.  The word $ab$ would also
be accepted, because the missing $b$-transition from the accepting state
$q_1$ would keep the automaton accepting.  This is not the intended
meaning of an omitted transition.
\end{example}

\begin{figure}[ht]
\centering
\begin{tikzpicture}[
  node distance=2.8cm,
  state/.style={circle, draw, thick, minimum size=11mm, fill=blue!10},
  accepting/.style={circle, draw, double, thick, minimum size=11mm, fill=teal!10},
  sink/.style={circle, draw, thick, minimum size=11mm, fill=gray!12},
  arr/.style={-Stealth, thick},
]
  \node[state] (q0) {$q_0$};
  \node[accepting, right=of q0] (q1) {$q_1$};
  \node[sink, right=of q1] (bot) {$\bot$};

  \draw[arr] ([xshift=-1.0cm]q0.west) -- (q0);
  \draw[arr] (q0) -- node[above] {$a$} (q1);
  \draw[arr] (q0) edge[bend right=22] node[below] {$b$} (bot);
  \draw[arr] (q1) -- node[above] {$a,b$} (bot);
  \draw[arr] (bot) edge[loop above] node {$a,b$} (bot);
\end{tikzpicture}
\caption{A concrete sink completion: every missing transition goes to the rejecting state $\bot$, and $\bot$ loops on every input symbol.}
\label{fig:concrete-sink-completion}
\end{figure}

\Cref{fig:concrete-sink-completion} shows the completed version of the
partial automaton from \cref{ex:missing-transition-not-self-loop}.  The
sink state means ``we have already violated the pattern''.  Once the
automaton enters it, no future input can make the word acceptable again.
This completion matters later, because complementation of deterministic
automata works only when the automaton is complete.


% ============================================================
\section{Non-Deterministic Finite Automata}
\label{sec:nfa}
% ============================================================

Non-determinism changes the transition structure.  The automaton may
have several possible next states for the same state and input symbol.
For example, while scanning a word, an automaton might keep reading as
usual, or it might guess that the current symbol is the important one.
The formal definition below records all such possible moves instead of
choosing one immediately.

\begin{definition}[Non-deterministic finite automaton]
\label{def:nfa}
A \textbf{non-deterministic finite automaton} (NFA) is a tuple
\[
  \mathcal{A}=(Q,A,\Delta,q_0,F)
\]
where $Q$, $A$, $q_0$, and $F$ are as for a DFA, and
\[
  \Delta \subseteq Q\times A\times Q
\]
is a transition relation.
\end{definition}

\aside{A DFA is just the special case where every pair $(q,a)$ has at
most one outgoing transition.  In an NFA, the same pair may lead to
several possible next states.}

Equivalently, we can view the transition structure as a function
\[
  \Delta:Q\times A\to 2^Q.
\]
The value $\Delta(q,a)$ is the set of all states that the automaton may
move to after reading $a$ in state $q$.

For computations it is convenient to extend this to sets of states.  If
$S\subseteq Q$, define
\[
  \widehat{\Delta}(S,\varepsilon)=S,
  \qquad
  \widehat{\Delta}(S,wa)=
    \bigcup_{p\in\widehat{\Delta}(S,w)}\Delta(p,a).
\]
Starting from a single state $q$, we write
$\widehat{\Delta}(q,w)$ for
$\widehat{\Delta}(\{q\},w)$.  The set
$\widehat{\Delta}(q_0,w)$ contains exactly the states reachable by some
run on $w$.

\begin{example}[Computing reachable states]
\label{ex:nfa-reachable-states}
Consider the NFA over $A=\{a,b\}$ shown in
\cref{fig:nfa-reachable-states}.  Its initial state is $q$, and we use
it only to compute reachability, so no accepting states are marked.
\end{example}

\begin{figure}[ht]
\centering
\begin{tikzpicture}[
  node distance=2.1cm,
  state/.style={circle, draw, thick, minimum size=10mm, fill=blue!10},
  arr/.style={-Stealth, thick},
]
  \node[state] (q) {$q$};
  \node[state, above right=0.7cm and 1.7cm of q] (q2) {$q_2$};
  \node[state, right=2.2cm of q2] (q1) {$q_1$};
  \node[state, below right=0.7cm and 1.7cm of q] (q4) {$q_4$};
  \node[state, right=2.2cm of q4] (q5) {$q_5$};
  \node[state, below left=0.7cm and 0.3cm of q5] (q3) {$q_3$};

  \draw[arr] ([xshift=-0.8cm]q.west) -- (q);
  \draw[arr] (q) to node[above left] {$b$} (q2);
  \draw[arr] (q) to node[below left] {$b$} (q4);
  \draw[arr] (q2) to node[above] {$b$} (q1);
  \draw[arr] (q1) to[bend left=12] node[above] {$a$} (q4);
  \draw[arr] (q4) to node[above] {$b$} (q5);
  \draw[arr, bend right=34] (q4) to node[below] {$b$} (q);
  \draw[arr] (q4) to[bend left=12] node[right] {$a$} (q1);
  \draw[arr] (q5) to node[right] {$a$} (q1);
  \draw[arr] (q5) to node[left] {$a$} (q3);
\end{tikzpicture}
\caption{A small NFA used to compute $\widehat{\Delta}(q,bba)$ step by step.}
\label{fig:nfa-reachable-states}
\end{figure}

Let us compute $\widehat{\Delta}(q,bba)$.
After the first $b$,
\[
  \widehat{\Delta}(q,b)=\{q_2,q_4\}.
\]
From $q_2$ another $b$ reaches $q_1$, while from $q_4$ another $b$
reaches both $q$ and $q_5$.  Hence
\[
  \widehat{\Delta}(q,bb)=\{q,q_1,q_5\}.
\]
The final $a$ kills the branch currently in $q$, because $q$ has no
outgoing $a$-transition.  The branch in $q_1$ reaches $q_4$, and the
branch in $q_5$ reaches both $q_1$ and $q_3$.  Therefore
\[
  \widehat{\Delta}(q,bba)=\{q_1,q_3,q_4\}.
\]

\aside{This is the set-based view of non-determinism: we do not choose
one branch and forget the others.  At each prefix, we collect all states
that are reachable by at least one run.}

\begin{definition}[Run and acceptance for an NFA]
\label{def:nfa-run}
Let $w=a_0a_1\cdots a_{n-1}$.  A \textbf{run} of an NFA on $w$ is a
sequence of states
\[
  r_0,r_1,\ldots,r_n
\]
such that $r_0=q_0$ and $(r_i,a_i,r_{i+1})\in\Delta$ for every
$0\leq i<n$.

The word $w$ is accepted if there exists at least one run on $w$ with
$r_n\in F$.
\end{definition}

The key word is \emph{exists}.  An NFA may have many failing runs on the
same input.  The word is accepted as soon as one run reaches an accepting
state after consuming the whole word.
\aside{Other acceptance choices are possible in principle, but the
standard NFA semantics is existential: one successful computation is
enough.}
Equivalently,
\[
  w\in L(\mathcal{A})
  \quad\Longleftrightarrow\quad
  \widehat{\Delta}(q_0,w)\cap F\neq\emptyset.
\]

\begin{example}[Guessing the last symbol]
\label{ex:nfa-last-a}
Over $A=\{a,b\}$, the language of words whose last symbol is $a$
can be recognized by a small NFA that stays in a scanning state, and
whenever it reads $a$, guesses that this might be the last symbol.
That guess moves to an accepting state with no outgoing transitions.
\end{example}

This example illustrates why non-determinism is convenient: the automaton
does not need to know in advance which $a$ is the last one.  It can try
all possibilities at once.

\begin{figure}[ht]
\centering
\begin{tikzpicture}[
  node distance=3.0cm,
  state/.style={circle, draw, thick, minimum size=11mm, fill=blue!10},
  accepting/.style={circle, draw, double, thick, minimum size=11mm, fill=teal!10},
  arr/.style={-Stealth, thick},
]
  \node[state] (scan) {$q_s$};
  \node[accepting, right=of scan] (last) {$q_a$};
  \draw[arr] ([xshift=-1.0cm]scan.west) -- (scan);
  \draw[arr] (scan) edge[loop above] node {$a,b$} (scan);
  \draw[arr, bend left=12] (scan) to node[above] {$a$} (last);
\end{tikzpicture}
\caption{An NFA for words ending in $a$.  The transition to $q_a$ is the guess that the current $a$ is the last symbol.}
\label{fig:nfa-last-a}
\end{figure}

Notice the visual difference between \cref{fig:even-a-dfa} and
\cref{fig:nfa-last-a}.  The DFA has exactly one outgoing transition for
each symbol at each state.  The NFA has two possible $a$-moves from
$q_s$: keep scanning, or guess that this $a$ is final.  If you draw a
second outgoing arrow labeled $a$ from the same state to another state,
you are already outside the DFA world, because one state-symbol pair
would have two possible targets.


% ============================================================
\section{Determinizing an NFA}
\label{sec:subset-construction}
% ============================================================

Every DFA is trivially an NFA: if the transition function $\delta$ is
total, we can view it as a relation $\Delta$ in which each pair $(q,a)$
has exactly one successor.  So the class of languages recognizable by
DFAs is contained in the class recognizable by NFAs.

The interesting question is the converse.  Are there languages that an
NFA can recognize but no DFA can?  Does non-determinism genuinely add
expressive power, or does it only add design convenience?

The answer is that it adds only convenience.  If we restrict ourselves to
DFAs, we do not lose any language.

\begin{theorem}[Subset construction]
\label{thm:subset-construction}
For every NFA $\mathcal{A}$ there exists a DFA $\mathcal{D}$ such that
\[
  L(\mathcal{D})=L(\mathcal{A}).
\]
\end{theorem}

\begin{formaldetails}[Construction]
Let $\mathcal{A}=(Q,A,\Delta,q_0,F)$ be an NFA.  Define the DFA
\[
  \mathcal{D}=(2^Q,A,\delta_D,\{q_0\},F_D)
\]
as follows:
\[
  \delta_D(S,a)=\bigcup_{q\in S}\Delta(q,a),
  \qquad
  F_D=\{S\subseteq Q:S\cap F\neq\emptyset\}.
\]
The deterministic state $S$ records the whole set of NFA states that can
be active after reading the current input prefix.
\end{formaldetails}

The construction is called the \textbf{power-set construction} because
the DFA states are subsets of the NFA state set.  Each deterministic
state $S$ is a subset of $Q$, so the DFA has at most $2^{|Q|}$ states.
The intuition is direct: for every possible subset of NFA states that
could be simultaneously active, the DFA introduces one dedicated state.
There is a price to pay in size, but not in expressive power --- the two
classes of automata define exactly the same set of languages.

\aside{This equivalence between DFA and NFA is specific to finite words.
When we move to automata over infinite words in
\cref{ch:omega-automata}, we will see that deterministic and
non-deterministic variants are \emph{not} always equivalent: the
acceptance condition matters, and for some conditions non-determinism
genuinely adds expressive power.}

\begin{figure}[ht]
\centering
\begin{tikzpicture}[
  setbox/.style={draw, rounded corners=4pt, fill=teal!8,
                 minimum width=23mm, minimum height=10mm, align=center},
  arr/.style={-Stealth, thick},
]
  \node[setbox] (s0) {$\{q_0\}$};
  \node[setbox, right=2.3cm of s0] (s1) {$\{q_0,q_a\}$};
  \node[setbox, right=2.3cm of s1] (s2) {$\emptyset$};
  \draw[arr] (s0) -- node[above] {$a$} (s1);
  \draw[arr] (s1) edge[loop above] node {$a$} (s1);
  \draw[arr] (s0) edge[loop above] node {$b$} (s0);
  \draw[arr, dashed, gray] (s2) edge[loop above] node {$a,b$} (s2);
\end{tikzpicture}
\caption{A determinized state is a set of currently possible NFA states.}
\label{fig:subset-state}
\end{figure}

\Cref{fig:subset-state} visualizes the key bookkeeping move: one
deterministic state stands for a whole set of possible NFA states.  The
picture is deliberately about sets, not individual branches.

\begin{proof}
We prove the invariant: after reading any word $w$, the state reached by
the DFA is exactly the set of states reachable by some NFA run on $w$.
For $w=\varepsilon$, both sides are $\{q_0\}$.  If the invariant holds
for $w$, then after reading $wa$ the DFA applies
\[
  \delta_D(S,a)=\bigcup_{q\in S}\Delta(q,a),
\]
which is precisely the set of all NFA successors reached by extending
one valid run on $w$ with one $a$-transition.  Thus the invariant holds
by induction on the word length.

The NFA accepts $w$ exactly when the reachable set contains some final
state, which is exactly the definition of $F_D$.
\end{proof}

% ============================================================
\subsection{Regular Languages}
\label{subsec:regular-languages}
% ============================================================

The equivalence between DFA and NFA lets us name the class of languages
that finite automata recognize.

\begin{definition}[Regular language]
\label{def:regular-language}
A language $L \subseteq A^*$ is \textbf{regular} if there exists a
DFA $\mathcal{A}$ with $L(\mathcal{A})=L$.
\end{definition}

This definition names the languages that can be checked with finite
memory.  To read it, ask whether there is some finite automaton whose
accept/reject behaviour matches the property exactly.  The operational
test is: build a DFA, run it on the whole input word, and accept exactly
when the final state lies in its accepting set.

\cref{thm:subset-construction} makes this definition robust to the
choice of model: a language is regular if and only if it is recognized
by some NFA.  We may therefore use DFA and NFA interchangeably when
reasoning about regularity --- typically NFAs when designing an
automaton, DFAs when we need a single run on each input, for example to
complement the language by swapping accepting and non-accepting states.

The class of regular languages is a \emph{strict} subset of all
languages over $A$.  This boundary is exactly what the theory is
about: regular languages are the properties that a finite agent can
check by reading the input once, using only a bounded amount of memory.
The next example shows a natural language that slips through this net.

\begin{example}[A non-regular language: $a^n b^n$]
\label{ex:anbn-non-regular}
Over $A=\{a,b\}$, consider
\[
  L_{=}=\{a^n b^n : n\geq 0\},
\]
the words consisting of a block of $a$'s followed by a block with the
same number of $b$'s.  This is precisely the matching problem that
\cref{ex:a-eventually-followed-by-b} deliberately avoided: there, a
single $b$ could discharge any number of pending $a$'s, whereas here
each $a$ demands a \emph{distinct} later $b$.

Suppose a DFA with $k$ states reads the unary block $a^i$ for
$i=0,1,\ldots,k$.  By the pigeonhole principle two of these prefixes,
say $a^i$ and $a^j$ with $i<j$, must lead to the same state.  From that
shared state, the same continuation $b^i$ must produce the same verdict
for both $a^i b^i$ and $a^j b^i$.  But $a^i b^i\in L_{=}$ while
$a^j b^i\notin L_{=}$, a contradiction.  Hence no DFA recognizes
$L_{=}$, and $L_{=}$ is not regular.

We will turn this pigeonhole sketch into a clean characterization ---
the Myhill--Nerode theorem --- in \cref{ch:regular-decision}.
\end{example}

\section*{Notation Introduced in This Chapter}
\addcontentsline{toc}{section}{Notation Introduced in This Chapter}

\begin{center}
\begin{tabular}{@{}ll@{}}
\toprule
\textbf{Notation} & \textbf{Meaning} \\
\midrule
$A$ & finite alphabet of input symbols \\
$A^*$ & set of all finite words over $A$ \\
$\varepsilon$ & empty word \\
$L \subseteq A^*$ & language over $A$ \\
$uv$ & concatenation of words $u$ and $v$ \\
$|w|$ & length of the word $w$ \\
$W^*$ & Kleene closure of the language $W$ \\
$W^+$ & non-empty finite concatenations of words from $W$ \\
$\mathcal{A}$ & finite automaton \\
$Q$ & finite set of states \\
$q_0$ & initial state \\
$F$ & set of accepting states \\
$\delta:Q\times A\to Q$ & DFA transition function \\
$\widehat{\delta}:Q\times A^*\to Q$ & extended DFA transition function on words \\
$L(\mathcal{A})$ & language accepted by the automaton $\mathcal{A}$ \\
$\bot$ & rejecting sink state \\
$\Delta \subseteq Q\times A\times Q$ & NFA transition relation \\
$\Delta:Q\times A\to 2^Q$ & equivalent set-valued view of NFA transitions \\
$\widehat{\Delta}$ & extended NFA transition map on words \\
$2^Q$ & powerset of $Q$, used as the state space in determinization \\
\bottomrule
\end{tabular}
\end{center}

\begin{summarybox}
\begin{itemize}[nosep]
  \item A language is a set of finite words over an alphabet.
  \item A DFA has one run per input word.
  \item An NFA may have many runs, and accepts if at least one succeeds.
  \item Every NFA can be determinized by tracking sets of possible states.
  \item The price of determinization can be exponential.
  \item The languages recognized by finite automata (DFA or NFA) are the
    \emph{regular} languages; not every language is regular, as
    $\{a^n b^n:n\geq 0\}$ shows.
\end{itemize}
\end{summarybox}

\section*{Exercises}
\addcontentsline{toc}{section}{Exercises}

\begin{chapex}[Simulating a DFA]
For the DFA in \cref{fig:even-a-dfa}, compute the reached state after
reading each prefix of the word $abbaab$.  Is the whole word accepted?
\end{chapex}

\begin{chapex}[One subset-construction step]
Use the NFA in \cref{fig:nfa-last-a}.  Starting from the determinized
state $\{q_s,q_a\}$, compute the next determinized state on input $b$.
Explain why the accepting state disappears.
\end{chapex}

\begin{chapex}[Design]
Design a DFA over $\{a,b\}$ for the language of words that contain the
substring $ab$.  Use a sink-like state to record that the substring has
already been seen.
\end{chapex}
